%% NrityaAI: A Real-Time Indian Classical Dance Style Classification
%% and Pose Correction System Using MediaPipe BlazePose and
%% Spatio-Temporal Deep Learning
%%
%% IEEE Access Draft
%% Author: Anoushka Dwivedi
%% Guide: Dr. Soumya A
%% RV College of Engineering, Bengaluru, Karnataka, India

\documentclass[journal]{IEEEtran}

\usepackage{cite}
\usepackage{amsmath,amssymb,amsfonts}
\usepackage{algorithmic}
\usepackage{graphicx}
\usepackage{textcomp}
\usepackage{xcolor}
\usepackage{booktabs}
\usepackage{hyperref}
\usepackage{array}
\usepackage{multirow}

\hypersetup{
    colorlinks=true,
    linkcolor=blue,
    citecolor=blue,
    urlcolor=blue
}

\begin{document}

\title{NrityaAI: A Real-Time Indian Classical Dance Style Classification
and Pose Correction System Using MediaPipe BlazePose and
Spatio-Temporal Deep Learning}

\author{Anoushka~Dwivedi,~\IEEEmembership{}
        and~Soumya~A,~\IEEEmembership{}%
\thanks{A. Dwivedi is with the Department of Computer Science and
Engineering (Data Science), RV College of Engineering, Bengaluru,
Karnataka 560059, India (e-mail: danoushka29@gmail.com).}%
\thanks{Dr. Soumya A is Professor and Programme Coordinator, CS-CD,
RV College of Engineering, Bengaluru, Karnataka 560059, India.}%
\thanks{Manuscript received May 2026.}}

\markboth{IEEE Access,~Vol.~XX, No.~XX, 2026}%
{Dwivedi and Soumya: NrityaAI}

\maketitle

% -----------------------------------------------------------------------
%  ABSTRACT
% -----------------------------------------------------------------------
\begin{abstract}
Indian Classical Dance (ICD) constitutes a living repository of
millennia-old cultural expression, yet the declining availability of
qualified instructors and the high cost of traditional tutelage have
created a critical accessibility gap that affects an estimated
300 million potential learners across the Indian subcontinent.
Existing computational approaches to dance recognition either address
a single dance form in isolation, operate exclusively on static images,
restrict analysis to hand-gesture recognition, or provide no real-time
corrective feedback—leaving the intersection of multi-form ICD
classification and intelligent pose correction entirely unaddressed.
This paper presents NrityaAI, the first unified, real-time system that
simultaneously classifies and provides pose correction feedback for
three canonical Indian Classical Dance forms: Bharatanatyam, Kathak,
and Odissi. The system employs MediaPipe BlazePose for
skeleton-based keypoint extraction, yielding 33 three-dimensional
landmarks per frame at ten frames per second, which are subsequently
processed by a novel dual-output Convolutional Neural Network and Long
Short-Term Memory (CNN+LSTM) spatio-temporal architecture. The network
produces two simultaneous outputs: a three-class style classification
head and a continuous pose quality scoring head. A YouTube-scraped
dataset of 51 source performance videos was processed into
keypoint-sequence windows, augmented with two Kaggle benchmark
collections to yield 8,329 labelled test sequences prior to
oversampling. The proposed model achieves 99.06\% overall classification
accuracy and a 99.37\% macro F1-score across the test partition,
with per-class F1-scores of 99.15\% for Bharatanatyam and 98.96\%
for Kathak. The Odissi evaluation is subject to a noted single-sample
limitation arising from systematic MediaPipe landmark occlusion in
source footage, which is acknowledged explicitly as a boundary condition
of the current study. NrityaAI constitutes a threefold contribution:
the first unified real-time classification system spanning these three
dance forms, a novel joint-angle deviation pose correction engine, and
the first publicly reproducible keypoint-sequence dataset for
Bharatanatyam, Kathak, and Odissi collectively. This work is aligned
with UNESCO Sustainable Development Goals SDG~4 (Quality Education),
SDG~11 (Sustainable Cities and Communities), and SDG~17 (Partnerships
for the Goals), advancing equitable access to intangible cultural
heritage through AI-driven pedagogy.
\end{abstract}

% -----------------------------------------------------------------------
%  INDEX TERMS
% -----------------------------------------------------------------------
\begin{IEEEkeywords}
Indian classical dance recognition, pose estimation, MediaPipe
BlazePose, convolutional neural networks, long short-term memory,
spatio-temporal deep learning, dance style classification, real-time
pose correction, cultural heritage preservation, skeleton-based action
recognition
\end{IEEEkeywords}

% -----------------------------------------------------------------------
%  SECTION I — INTRODUCTION
% -----------------------------------------------------------------------
\section{Introduction}
\label{sec:introduction}

Indian Classical Dance (ICD) represents one of the most sophisticated
and codified performing arts traditions in human history. Bharatanatyam,
originating in the temples of Tamil Nadu and formalised through the
\textit{Natya Shastra}—a treatise on performance arts attributed to
the sage Bharata Muni and dated between 200~BCE and 200~CE—is
characterised by its angular, rhythmic footwork, expressive
\textit{mudras} (hand gestures), and intricate \textit{abhinaya}
(emotive expression). Kathak, which evolved through the courts of
Mughal and Rajput patronage in northern India, is distinguished by its
rapid spins (\textit{chakkar}), syncopated footwork, and fluid upper
body vocabulary. Odissi, indigenous to the temples of Odisha, is
marked by its distinctive tribhanga (three-bend) posture, lyrical
quality, and sculptural aesthetic rooted in temple iconography.
All three forms have received recognition from cultural bodies
including UNESCO and the Sangeet Natak Akademi, and collectively
represent an irreplaceable segment of India's intangible cultural
heritage. Despite their historical and cultural primacy, the
transmission of these forms remains critically dependent on the
traditional \textit{guru-shishya parampara}—the master-disciple
pedagogical lineage—which has experienced accelerating decline in
the twenty-first century.

The erosion of access to classical dance education is a multidimensional
problem. Qualified \textit{gurus} are concentrated in metropolitan
centres, creating a severe geographic imbalance that effectively
excludes learners in rural and semi-urban settings. The financial cost
of sustained instruction is prohibitive for large segments of the
population, particularly when compounded by the need for appropriate
performance attire, musical accompaniment, and dedicated practice
space. Importantly, even where instruction is nominally available
through video-sharing platforms such as YouTube, the learner receives
no interactive feedback; the medium remains passive and unidirectional,
unable to identify misaligned postures, incorrect joint angles, or
stylistic drift from canonical forms. These barriers collectively
deprive an estimated 300 million Indians—many of whom express a desire
to engage with their classical artistic heritage—of meaningful
educational access. The digital preservation of these art forms
acquires additional urgency in light of the broader cultural risk:
as practising \textit{gurus} age and the apprentice pool contracts,
kinematic and expressive knowledge that is insufficiently documented
risks irreversible loss.

Artificial intelligence and computer vision offer a meaningful pathway
toward democratising access to classical dance education. A system
capable of real-time pose estimation and classification could serve as
a virtual teaching assistant, delivering the kind of corrective feedback
previously available only in direct pedagogical settings. Such a system
would be particularly transformative if it operated on commodity
hardware—consumer-grade laptops and webcams—without requiring
specialised depth sensors or high-performance graphical processing
units.

A careful survey of the existing literature reveals that prior
computational work on ICD is fragmented along several axes. Static
image-based classifiers~\cite{jain2021} achieve competitive accuracy on
benchmark datasets but are fundamentally incapable of capturing the
temporal dynamics—rhythmic patterns, transitional movements, and
temporal phrasing—that constitute the grammar of these dance forms.
Systems that address temporal structure typically focus on a single
dance form, precluding the cross-form discrimination that a unified
system requires~\cite{naik2021}. A distinct body of work addresses
hand-gesture or \textit{mudra} recognition in isolation~\cite{challapalli2022},
treating the upper extremity as the sole site of meaning while
disregarding the full-body kinematics that distinguish, for instance, a
Bharatanatyam \textit{aramandi} from a Kathak \textit{tatkaar}. A 2024
survey published in Elsevier's \textit{Entertainment Computing}
journal~\cite{survey2024} explicitly confirms that no unified real-time
system currently exists that jointly classifies and provides corrective
feedback for multiple ICD forms. Parallel work in the adjacent domain
of yoga pose correction~\cite{thoutam2022,swain2022} demonstrates the
feasibility of CNN+LSTM architectures for real-time skeletal feedback,
but these systems are developed exclusively for static or quasi-static
yoga postures and do not generalise to the highly dynamic temporal
vocabulary of classical dance.

NrityaAI addresses all of the foregoing gaps within a single, end-to-end
deployable framework. The system combines MediaPipe BlazePose's
real-time 33-landmark skeleton extraction with a novel dual-output
CNN+LSTM architecture that simultaneously produces a three-class
dance style prediction and a continuous pose quality score. A
joint-angle deviation engine computes angular error at six critical
skeletal joints and generates human-interpretable corrective cues,
enabling the learner to receive directional and magnitude-qualified
feedback in real time. The complete system is deployed as a web
application via a FastAPI backend and a Streamlit frontend, operational
on CPU-only consumer hardware.

The specific contributions of this work are as follows:
\begin{enumerate}
    \item The first unified, real-time classification system spanning
    Bharatanatyam, Kathak, and Odissi, achieving 99.06\% overall
    accuracy on a held-out test partition.
    \item A novel dual-output CNN+LSTM spatio-temporal architecture
    that jointly optimises for dance style classification and
    pose quality scoring in a single forward pass.
    \item A joint-angle deviation pose correction engine that quantifies
    angular error at six skeletal joints and generates per-joint
    directional feedback, operationalised as a real-time corrective
    overlay.
    \item The first unified keypoint-sequence dataset for Bharatanatyam,
    Kathak, and Odissi, constructed through systematic YouTube scraping
    with MediaPipe-based annotation and augmented with two Kaggle
    benchmark collections.
    \item A costume-invariant and background-invariant skeleton-based
    representational scheme that transcends the pixel-level feature
    limitations that affect RGB-image-based classifiers operating on
    visually diverse traditional performance attire.
\end{enumerate}

The remainder of this paper is structured as follows.
Section~\ref{sec:related} surveys related work across pose estimation
foundations, prior ICD recognition systems, and pose correction
frameworks. Section~\ref{sec:methodology} describes the complete
NrityaAI pipeline, including dataset construction, keypoint extraction,
feature engineering, model architecture, training configuration, the
pose correction engine, and system deployment. Section~\ref{sec:setup}
details the experimental setup, evaluation metrics, and hardware and
software environment. Section~V will present experimental results and
comparative analysis.

% -----------------------------------------------------------------------
%  SECTION II — RELATED WORK
% -----------------------------------------------------------------------
\section{Related Work}
\label{sec:related}

\subsection{Pose Estimation Foundations}

The computational estimation of human body pose from monocular image or
video input has advanced substantially over the past decade. The
landmark contribution of Cao et al.~\cite{openpose2017}—the OpenPose
framework—introduced Part Affinity Fields (PAFs), a set of learned
two-dimensional vector fields encoding the orientation and location of
body part connections. Operating on an 18-keypoint skeletal topology,
OpenPose achieves high-quality pose estimates across diverse scenes and
multiple concurrent subjects. However, its inference pipeline is
computationally intensive, requiring a dedicated GPU to achieve
frame-rates suitable for real-time applications; deployment on
consumer CPU hardware produces latencies that render the system
impractical for interactive pedagogy. Furthermore, the 18-keypoint
topology captures coarse limb positions but underrepresents the foot
and ankle articulation—\textit{pada} positions—that are biomechanically
central to both Bharatanatyam and Kathak footwork vocabulary.

Bazarevsky et al.~\cite{blazepose2020} addressed these limitations
with BlazePose, the pose estimation backbone of the MediaPipe
framework. BlazePose employs a two-stage pipeline consisting of a
lightweight pose detector followed by a regression-based landmark
localiser, yielding 33 anatomical landmarks with three-dimensional
coordinates and per-landmark visibility scores in a single
forward pass. The architecture is explicitly optimised for real-time
inference on mobile and CPU environments, achieving 30+ frames per
second (FPS) on commodity hardware without GPU acceleration.
Critically, BlazePose's 33-keypoint topology includes dedicated
landmarks for the heels, toe tips, and facial features, providing
substantially richer coverage of the kinematic degrees of freedom
relevant to classical dance. For the present work, BlazePose was
selected over OpenPose for three principal reasons: its
CPU-deployability on consumer hardware, its superior foot-region
landmark density, and its provision of landmark visibility scores
that permit principled filtering of occluded or unreliable keypoints.

\subsection{Indian Classical Dance Recognition}

Jain et al.~\cite{jain2021} presented a convolutional classification
system for eight Indian classical dance forms using a ResNet-50
backbone pre-trained on ImageNet and fine-tuned on a curated corpus
of performance images. The system achieved 91.1\% top-1 accuracy,
demonstrating the discriminative power of deep convolutional features
for inter-form style distinction. However, the architecture operates
on individual static frames and therefore cannot model the temporal
dynamics that constitute the most structurally distinctive features of
each dance form. A Bharatanatyam \textit{alarippu} sequence and an
Odissi \textit{mangalacharan} may contain individual frames with
superficially similar full-body configurations; it is the temporal
ordering and rhythmic patterning of those configurations that
unambiguously identifies the form. Static classifiers are structurally
incapable of leveraging this information.

Challapalli and Devarakonda~\cite{challapalli2022} proposed a hybrid
optimisation framework combining Convolutional Neural Networks with
Particle Swarm Optimisation (PSO) and Grey Wolf Optimisation (GWO)
for feature selection in dance pose classification. While the
meta-heuristic feature selection provides computational efficiency
gains, the fundamental representational substrate remains pixel-level
RGB features extracted from still images. This approach inherits all of
the limitations of image-based classifiers—sensitivity to background
clutter, illumination variation, and crucially, the high intra-class
visual variability introduced by the elaborate and stylistically
diverse traditional costumes worn across performances of the same
dance form.

Naik and Supriya~\cite{naik2021} explored an alternative
representational modality, applying the PointNet deep learning
architecture to three-dimensional point cloud representations of dance
poses. This skeleton-free volumetric approach avoids the costume
sensitivity of RGB classifiers and provides genuine three-dimensional
spatial information. However, it requires depth camera hardware—such
as Microsoft Kinect or Intel RealSense sensors—that is unavailable
in typical home or studio settings. The practical deployment barrier
associated with depth sensor requirements critically limits the
accessibility advantage that an AI-based dance teaching system is
intended to provide.

A comprehensive survey published in Elsevier's \textit{Entertainment
Computing} journal~\cite{survey2024} examined the state of the art in
ICD recognition and drew the explicit conclusion that no existing system
combines multi-form classification with real-time corrective feedback
in a deployable pipeline. The survey identified the absence of
large-scale, annotated ICD keypoint datasets as a primary impediment
to further progress. A contemporaneous study published in
IEEE Xplore~\cite{ieee2025} applied convolutional architectures
to dance pose classification and achieved competitive accuracy metrics,
but similarly produced no pose correction output, leaving the
feedback dimension unaddressed.

\subsection{Pose Correction Systems}

The closest technical analogue to the NrityaAI pose correction engine
is found in the domain of yoga posture assessment. Thoutam et
al.~\cite{thoutam2022} constructed a CNN+LSTM pipeline for yoga pose
classification and correctness assessment, achieving 99.38\% accuracy
across a set of canonical yoga asanas. The temporal model captures the
transition dynamics of sequential pose entry, and the system provides
binary correct/incorrect feedback per pose. While the architecture
pattern closely parallels the design decisions adopted in the present
work, the system addresses quasi-static yoga postures—poses that are
held for extended durations—and is not designed for the continuous,
rapidly evolving kinematic trajectories characteristic of classical
dance performance.

Swain et al.~\cite{swain2022} implemented a real-time yoga monitoring
system using MediaPipe for keypoint extraction and a CNN+LSTM
architecture for temporal modelling, targeting six canonical yoga
poses. The use of MediaPipe as the keypoint extraction backbone is
directly analogous to the present work, and the real-time operational
requirement is shared. The limitation is scope: with six yoga poses,
the inter-class kinematic diversity is substantially lower than the
multi-form, continuous-sequence discrimination required for ICD
classification. A 2024 IEEE study~\cite{yogavision2024} further
refined this class of system with a CNN+LSTM+MediaPipe fusion
architecture for yoga pose correction, achieving high accuracy in
the yoga domain but remaining inapplicable to dance forms.

\subsection{Research Gap Summary}

The foregoing survey reveals a clear and multidimensional research gap.
Existing work on Indian Classical Dance recognition is either restricted
to static image classification, addresses a single dance form in
isolation, or focuses narrowly on sub-components such as hand gestures.
Systems that address temporal dynamics in movement recognition are
confined to the yoga domain, operate on highly constrained pose
vocabularies, and provide binary rather than quantitative directional
feedback. No prior system jointly addresses multi-form ICD
classification, real-time temporal modelling of dance sequences, and
per-joint corrective feedback within a single deployable architecture.
NrityaAI directly targets this compound gap, providing the first unified
system that simultaneously handles classification across Bharatanatyam,
Kathak, and Odissi, delivers continuous pose quality scores, and
generates joint-specific directional correction cues—all on
CPU-only consumer hardware without specialised sensors.

% -----------------------------------------------------------------------
%  SECTION III — METHODOLOGY
% -----------------------------------------------------------------------
\section{Methodology}
\label{sec:methodology}

\subsection{System Overview}

NrityaAI is structured as a sequential processing pipeline that
transforms raw video input into dual classification and pose correction
outputs in real time. The pipeline proceeds through five principal
stages. In the first stage, source video is ingested either from a
file upload or from a live webcam stream via the Streamlit frontend.
In the second stage, each video frame is passed through the MediaPipe
BlazePose inference engine, which extracts a 33-landmark three-dimensional
skeleton and associates per-landmark visibility scores. In the third
stage, the extracted keypoint sequences undergo coordinate normalisation
and are assembled into fixed-duration sliding windows that constitute
the model's input representation. In the fourth stage, the pretrained
dual-output CNN+LSTM model processes each window in a single forward
pass, producing a softmax class probability distribution over
\{Bharatanatyam, Kathak, Odissi\} and a sigmoid pose quality score in
the interval $[0, 1]$. In the fifth stage, a separate joint-angle
deviation engine computes angular errors at six critical skeletal
joints by comparing the live pose to a reference pose library, and
generates human-readable corrective instructions that are overlaid
on the skeleton visualisation. The complete pipeline is served by a
FastAPI backend and rendered through a Streamlit frontend, forming an
end-to-end web application capable of operating at target inference
speeds exceeding 20 FPS on CPU-only hardware. A schematic of the
full pipeline is shown in Fig.~\ref{fig:architecture}.

\begin{figure}[!t]
\centering
\includegraphics[width=\columnwidth]{figures/system_architecture}
\caption{NrityaAI system architecture. The pipeline transforms raw
video input through MediaPipe keypoint extraction, feature engineering,
CNN+LSTM inference, and a joint-angle correction engine, rendered via a
FastAPI–Streamlit web interface.}
\label{fig:architecture}
\end{figure}

\subsection{Dataset Construction}

\subsubsection{YouTube Scraping Strategy}

Primary training data was collected by systematically scraping
performance videos from YouTube using the \texttt{yt-dlp} command-line
utility. Search queries were constructed to retrieve high-quality,
professionally staged performances of each target dance form, prioritising
videos with clear full-body visibility, controlled stage lighting, and
minimal occlusion from co-performers or stage furniture. The final
corpus comprised 15 Bharatanatyam source videos, 19 Kathak source
videos, and 17 Odissi source videos, for a total of 51 source recordings.
Each video was downloaded at its native resolution and subsequently
processed by the MediaPipe extraction pipeline. To prevent temporal
data leakage between training, validation, and test partitions, the
split was performed at the video level rather than the frame or window
level; frames from any single source video appear in exactly one of the
three partitions.

\subsubsection{Supplementary Kaggle Sources}

Two publicly available Kaggle datasets were incorporated to supplement
the YouTube-scraped corpus. The HackerEarth Indian Dance Forms dataset
provides 599 images distributed across 8 classical dance style classes,
of which the three classes relevant to NrityaAI (Bharatanatyam, Kathak,
Odissi) were retained. The Bharatanatyam Dance Poses dataset provides
731 images across 9 pose categories and was used to enrich the
reference pose library employed by the pose correction engine. A
contextual transfer learning baseline was also established using the
AIST++ dataset—a large-scale professional dance dataset with 3D
annotations—to calibrate the temporal modelling capacity of the
CNN+LSTM architecture prior to fine-tuning on ICD-specific data.

\subsubsection{Class Distribution and Oversampling}

MediaPipe BlazePose landmark extraction was applied to all frames of all
51 source videos. The extraction produced a markedly imbalanced class
distribution: Bharatanatyam yielded 4,553 sliding-window sequences,
Kathak yielded 3,775 sequences, and Odissi yielded a single valid
window. The near-total collapse of Odissi representation is attributable
to systematic landmark occlusion in the Odissi source footage: the
characteristic \textit{tribhanga} posture, combined with heavy
traditional jewellery and the frequent use of low-key stage lighting
in Odissi productions, caused MediaPipe's landmark visibility scores
to fall below the acceptance threshold of $v \geq 0.5$ on the majority
of frames, rendering those frames unusable. This limitation is
explicitly acknowledged as a boundary condition of the current system;
Odissi results must be interpreted with the awareness that the evaluation
is effectively performed on a single test instance, and performance
claims for Odissi do not carry the same statistical weight as those
for Bharatanatyam and Kathak.

To address the resulting class imbalance, a random oversampling strategy
was applied to the training partition: sequences from underrepresented
classes were sampled with replacement until all three classes reached
the cardinality of the largest class (4,553 sequences per class),
yielding 13,659 training samples post-balancing. The test partition was
not oversampled; evaluation was performed on the unmodified distribution
of 8,329 test sequences to preserve honest assessment of real-world
performance.

\subsection{Keypoint Extraction}

The MediaPipe BlazePose inference engine was configured with the
following parameters: \texttt{static\_image\_mode=False},
\texttt{model\_complexity=1}, \texttt{min\_detection\_confidence=0.5},
and \texttt{min\_tracking\_confidence=0.5}. Model complexity level 1
represents the medium-accuracy variant, which provides a
favourable balance between landmark localisation precision and
real-time inference throughput. Source video was sampled at 10 FPS
prior to landmark extraction; this sampling rate was chosen to capture
the relevant kinematic dynamics of classical dance movements—whose
characteristic tempos (\textit{laya}) range from slow (\textit{vilambita
laya}) to fast (\textit{drut laya})—while limiting computational cost
and dataset size. For each accepted frame, BlazePose yields 33 landmark
records, each comprising a normalised x-coordinate, a normalised
y-coordinate, a metric-scale z-coordinate (depth relative to the hip
midpoint), and a scalar visibility score $v \in [0, 1]$. Frames for
which more than 25\% of landmarks registered $v < 0.5$ were discarded
as insufficiently reliable; the visibility threshold ensures that
partially occluded or heavily motion-blurred frames do not corrupt
the temporal windows. The output of the extraction stage is stored as
CSV files with 134 columns per row: one frame identifier, 33 groups
of four features (x, y, z, v), and a class label.

\subsection{Feature Engineering}

\subsubsection{Coordinate Normalisation}

Raw BlazePose coordinates are expressed relative to the image frame
and therefore encode camera distance and framing as confounds. To
produce a representation invariant to scale and camera placement, each
frame's landmark set is normalised relative to the performer's own body.
The reference origin is set to the midpoint of the left and right hip
landmarks (MediaPipe indices 23 and 24). A scale factor is derived from
the torso height, defined as the Euclidean distance from the hip
midpoint to the midpoint of the left and right shoulder landmarks
(indices 11 and 12). All landmark x, y, and z coordinates are
translated to the hip midpoint origin and then divided by the torso
height scale factor, mapping coordinates into the interval $[-1, 1]$
for typical full-body poses. This normalisation scheme ensures that
the feature representation is invariant to performer stature,
camera zoom level, and spatial position within the frame.

\subsubsection{Temporal Windowing}

Normalised landmark sequences are segmented into fixed-duration windows
using a sliding window approach with a window length of 60 frames and
a step size of 30 frames, corresponding to 50\% overlap between
consecutive windows. At the 10 FPS sampling rate, each window
represents 6 seconds of dance performance—a duration chosen to encompass
a complete rhythmic cycle (\textit{avartana}) at moderate tempo for all
three dance forms. The 50\% overlap ensures that no meaningful movement
phrase falls exclusively at the boundary of a single window. The output
shape of each sample fed to the model is $(60, 33, 4)$: 60 temporal
frames, 33 landmarks per frame, and 4 features (x, y, z, v) per
landmark.

\subsubsection{Joint Angle Computation}

To supplement the raw coordinate features with biomechanically
interpretable quantities, joint angles are computed at six skeletal
joints using the \texttt{arctan2} formulation of the two-bone angle.
For a three-landmark chain $(A, B, C)$ representing the proximal
segment, joint, and distal segment respectively, the angle $\theta$ at
joint $B$ is computed as:
\begin{equation}
    \theta = \left| \text{arctan2}\left( C_y - B_y,\; C_x - B_x \right)
    - \text{arctan2}\left( A_y - B_y,\; A_x - B_x \right) \right|
    \label{eq:angle}
\end{equation}

The six joints instrumented are: left elbow (landmarks 11, 13, 15),
right elbow (12, 14, 16), left knee (23, 25, 27), right knee (24, 26, 28),
left hip (11, 23, 25), and right hip (12, 24, 26). Joint angles are
used by the pose correction engine described in
Section~\ref{subsec:correction} but are not included in the CNN+LSTM
input tensor, which operates directly on the normalised coordinate
features.

\subsubsection{Data Partitioning}

Dataset partitioning was performed at the source video level to prevent
temporal data leakage. Source videos were randomly assigned to training,
validation, and test sets at a 70/15/15 ratio. All sliding windows
extracted from a given source video appear in exactly one partition;
no frame from a training video appears in the validation or test
partitions. This video-level split ensures that the model is evaluated
on genuinely unseen performance instances rather than on temporally
adjacent windows from the same recording, providing a rigorous assessment
of generalisation capacity.

\subsection{Model Architecture}
\label{subsec:architecture}

The NrityaAI model is a dual-output spatio-temporal neural network
designed to simultaneously classify dance style and estimate pose
quality from a sequence of full-body keypoint frames.
Fig.~\ref{fig:model} illustrates the complete architecture.

\begin{figure}[!t]
\centering
\includegraphics[width=\columnwidth]{figures/model_architecture}
\caption{NrityaAI dual-output CNN+LSTM architecture. Spatial features
are extracted per frame by TimeDistributed convolutional layers,
temporal dependencies are modelled by stacked LSTM layers, and two
independent dense heads produce the style classification and pose
quality outputs.}
\label{fig:model}
\end{figure}

The input tensor has shape $(\text{batch}, 60, 33, 4)$, representing
a batch of 60-frame windows, each frame containing 33 landmarks with
4 features each. The 33-landmark representation is treated as a
one-dimensional spatial sequence of length 33 over 4 channels,
appropriate for one-dimensional convolution.

\subsubsection{Spatial Feature Extraction}

Spatial features are extracted using a pair of
\texttt{TimeDistributed} one-dimensional convolutional layers.
The \texttt{TimeDistributed} wrapper applies an identical convolutional
encoder independently to each of the 60 temporal frames, enforcing
parameter sharing across time and ensuring that the same spatial
feature extractor operates on every frame in the sequence. The first
layer applies 64 filters of kernel size 3 with ReLU activation and
same-padding; the second applies 128 filters of kernel size 3 with
ReLU activation and same-padding. A \texttt{TimeDistributed
GlobalAveragePooling1D} layer then aggregates the spatial
feature maps into a fixed-length descriptor per frame. The output
of this stage has shape $(\text{batch}, 60, 128)$: a 128-dimensional
spatial feature vector for each of the 60 temporal frames. The use
of convolutional layers for spatial feature extraction is motivated
by the local anatomical structure of the skeleton: adjacent landmarks
in the BlazePose topology correspond to anatomically proximate body
parts, and the convolutional receptive field is well-suited to
capturing local joint configurations such as elbow flexion or
knee angulation.

\subsubsection{Temporal Modelling}

Temporal dependencies across the 60-frame sequence are modelled by a
stack of two Long Short-Term Memory (LSTM) layers. The first LSTM
layer contains 256 units and is configured with
\texttt{return\_sequences=True}, outputting a hidden state at every
time step to preserve temporal detail for the subsequent LSTM layer.
A Dropout layer with rate 0.3 follows the first LSTM to mitigate
overfitting to idiosyncratic temporal patterns in the training corpus.
The second LSTM layer contains 128 units and operates in
many-to-one mode (\texttt{return\_sequences=False}), consuming the
full sequence and emitting a single 128-dimensional context vector
that summarises the temporal dynamics of the 6-second window.
A second Dropout layer with rate 0.3 follows. LSTM was selected over
Gated Recurrent Units (GRU) for the temporal modelling stage because
the 60-frame sequence length demands robust long-range dependency
capture; LSTM's separate cell state and hidden state mechanisms
provide a more expressive gradient pathway for retaining information
from early frames while processing later frames compared to GRU's
unified state formulation.

\subsubsection{Dual Output Heads}

The 128-dimensional context vector produced by the second LSTM layer
is consumed by two independent dense output heads in parallel.
The style classification head consists of a single \texttt{Dense}
layer with 3 units and a softmax activation, producing a probability
distribution over the three dance classes:
$\hat{y}_\text{style} \in \mathbb{R}^3$ with
$\sum_k \hat{y}_k = 1$. The pose quality scoring head consists of
a single \texttt{Dense} layer with 1 unit and a sigmoid activation,
producing a scalar pose quality estimate $\hat{y}_\text{score}
\in [0, 1]$. The dual-head design enables simultaneous optimisation
for both objectives in a single forward and backward pass, avoiding
the computational overhead of two independent inference pipelines.
The skeleton-based input representation—derived from normalised
landmark coordinates rather than raw pixel values—provides inherent
invariance to costume appearance, background content, and image-level
photometric variation, addressing a key failure mode of RGB-based
classifiers applied to the visually diverse world of classical
dance performance.

\subsection{Training Configuration}

The model was compiled with the Adam optimiser at a default learning
rate of $\eta = 0.001$. The total training loss was defined as a
weighted sum of two component objectives: categorical cross-entropy
on the style classification head with a weight of 1.0, and mean
squared error (MSE) on the pose quality scoring head with a weight
of 0.5. It must be noted that the pose quality scoring head was
trained on synthetically constructed labels (constant value 0.5 per
window, representing a neutral quality baseline) in the absence of
human-expert quality annotations for the training corpus. The
implications of this training condition are discussed in
Section~\ref{sec:setup} as an explicit limitation of the current
system.

Three training callbacks were employed. An EarlyStopping callback
monitored validation style accuracy with a patience of 10 epochs,
halting training when no improvement was observed for 10 consecutive
epochs. A ModelCheckpoint callback saved the parameter configuration
achieving the highest validation accuracy over the training run.
A ReduceLROnPlateau callback halved the learning rate (factor 0.5)
when validation loss failed to improve for 5 consecutive epochs,
enabling the optimiser to navigate the loss landscape with finer
granularity in later training phases. Training proceeded for 88 epochs
before the EarlyStopping criterion was met, with the best checkpoint
recorded at epoch 82. All training was conducted on CPU-only
Apple Silicon hardware, with Metal Performance Shaders (MPS) available
as a fallback for PyTorch-based operations. Training and validation
accuracy and loss curves are shown in Fig.~\ref{fig:training}.

\begin{figure}[!t]
\centering
\includegraphics[width=\columnwidth]{figures/training_history}
\caption{NrityaAI training history. (a) Style classification accuracy
on training and validation sets across 88 epochs. (b) Total training
loss across epochs, showing convergence under ReduceLROnPlateau
scheduling.}
\label{fig:training}
\end{figure}

\subsection{Pose Correction Engine}
\label{subsec:correction}

The pose correction engine operates in parallel with the classification
head, consuming the same MediaPipe keypoint stream but executing a
separate angular deviation analysis. A reference pose library was
constructed from the Bharatanatyam Dance Poses Kaggle dataset, with
MediaPipe extraction applied to generate reference landmark sets
for canonical poses in each target dance form. For each incoming
live frame, joint angles at the six instrumented joints (left and
right elbows, knees, and hips) are computed using
Equation~\ref{eq:angle} for both the live pose and the nearest
reference pose, as determined by minimum Euclidean distance in the
normalised coordinate space. The angular deviation at each joint
$j$ is defined as $\Delta\theta_j = |\theta_j^\text{live}
- \theta_j^\text{ref}|$. A joint is flagged for correction if
$\Delta\theta_j \geq 15^\circ$, a threshold calibrated empirically
to balance sensitivity to meaningful postural error against resilience
to minor natural kinematic variation. For each flagged joint, the
correction engine generates a human-readable feedback string
specifying the joint identity, the direction of required adjustment
(increase or decrease), and the magnitude of deviation in degrees.
An aggregate pose quality score is computed as:
\begin{equation}
    S_\text{pose} = \max\left(0,\; 100 - \sum_{j=1}^{6} \Delta\theta_j\right)
    \label{eq:score}
\end{equation}
This score is displayed to the learner as $S_\text{pose} \times 0.5$
on a 50-point scale in the frontend interface, providing an intuitively
bounded quality metric. The correction feedback and pose score are
rendered as a real-time overlay on the skeleton visualisation, drawn
using OpenCV's graphic primitives, and delivered to the learner with
sub-second latency.

\subsection{System Deployment}

NrityaAI is deployed as a two-tier web application. The backend
is implemented in FastAPI and exposes three REST endpoints:
\texttt{/analyze-video} for batch processing of uploaded video files,
\texttt{/analyze-frame} for real-time single-frame inference from
webcam streams, and \texttt{/health} for service status monitoring.
The frontend is implemented in Streamlit and provides two operational
modes: a video upload tab for post-hoc analysis of recorded performances,
and a live webcam tab for real-time interactive practice sessions.
In both modes, the output display includes the predicted dance style
label with confidence probability, the pose quality score, per-joint
correction cues, and a skeleton overlay rendered via OpenCV drawing
utilities on the original video frame. The system targets an inference
throughput exceeding 20 FPS at the \texttt{/analyze-frame} endpoint
on CPU-only consumer hardware, ensuring a real-time interactive
experience without hardware specialisation requirements.

% -----------------------------------------------------------------------
%  SECTION IV — EXPERIMENTAL SETUP
% -----------------------------------------------------------------------
\section{Experimental Setup}
\label{sec:setup}

\subsection{Hardware and Software Environment}

All experiments were conducted on a MacBook equipped with an Apple
Silicon processor. Model training and evaluation were performed
entirely on CPU, with Metal Performance Shaders (MPS) utilised as a
fallback accelerator for PyTorch-based auxiliary operations. No
dedicated GPU was employed at any stage of the experimental pipeline,
establishing that the NrityaAI system is trainable and deployable
on consumer-grade hardware.

The software stack comprised Python 3.10, TensorFlow/Keras 2.12+
for model definition and training, MediaPipe 0.10+ for landmark
extraction, FastAPI 0.100+ for API serving, Streamlit 1.25+ for the
frontend interface, OpenCV 4.8+ for video processing and skeleton
visualisation, scikit-learn for evaluation metric computation and
train/validation/test split management, NumPy for numerical array
operations, and Pandas for dataset management and CSV-format data
storage.

\subsection{Dataset Statistics}

Table~\ref{tab:dataset} summarises the dataset composition before and
after oversampling, and the test partition used for final evaluation.

\begin{table}[!t]
\caption{Dataset Statistics by Class}
\label{tab:dataset}
\centering
\begin{tabular}{lcccc}
\toprule
\textbf{Class} & \textbf{Source} & \textbf{Raw} & \textbf{Post-} & \textbf{Test}\\
& \textbf{Videos} & \textbf{Windows} & \textbf{Oversample} & \textbf{Samples}\\
\midrule
Bharatanatyam & 15 & 4,553 & 4,553 & 4,553 \\
Kathak        & 19 & 3,775 & 4,553 & 3,775 \\
Odissi$^{*}$  & 17 & 1     & 4,553 & 1 \\
\midrule
\textbf{Total} & \textbf{51} & \textbf{8,329} & \textbf{13,659} & \textbf{8,329} \\
\bottomrule
\multicolumn{5}{l}{$^{*}$Odissi oversampled from 1 real window; results for this}\\
\multicolumn{5}{l}{class must be interpreted with caution (see Section~\ref{subsec:limitations}).}\\
\end{tabular}
\end{table}

\subsection{Evaluation Metrics}

Model performance was assessed using the following metrics, computed
over the held-out test partition. Overall accuracy is defined as the
ratio of correctly classified test windows to total test windows.
Per-class precision, recall, and F1-score were computed according to
the standard binary one-versus-rest formulation for each of the three
classes. The macro-averaged F1-score—the unweighted mean of per-class
F1-scores—serves as the primary aggregate performance indicator, as
it treats each class equally and is sensitive to class-level failure
modes. The confusion matrix provides a complete accounting of
inter-class error patterns. For the pose correction engine, the
Mean Angle Error (MAE) across the six instrumented joints is computed
as a secondary metric, quantifying the fidelity of the angular
deviation estimates.

\subsection{Training Hyperparameters}

Table~\ref{tab:hyperparams} lists the complete hyperparameter
configuration used for the reported experiments.

\begin{table}[!t]
\caption{Training Hyperparameter Configuration}
\label{tab:hyperparams}
\centering
\begin{tabular}{lc}
\toprule
\textbf{Hyperparameter} & \textbf{Value} \\
\midrule
Window size & 60 frames \\
Step size & 30 frames \\
Keypoints per frame & 33 \\
Features per keypoint & 4 \\
CNN filters (layer 1 / layer 2) & 64 / 128 \\
LSTM units (layer 1 / layer 2) & 256 / 128 \\
Dropout rate & 0.3 \\
Batch size & 32 \\
Maximum epochs & 100 \\
Actual training epochs & 88 \\
Early stopping patience & 10 \\
Initial learning rate & 0.001 (Adam) \\
LR reduction factor & 0.5 \\
LR reduction patience & 5 \\
Style loss weight & 1.0 (CCE) \\
Pose score loss weight & 0.5 (MSE) \\
\bottomrule
\end{tabular}
\end{table}

\subsection{Acknowledged Limitations}
\label{subsec:limitations}

Two limitations of the experimental setup require explicit
acknowledgement. First, the Odissi class is represented by a single
valid test window, arising from the systematic occlusion described in
Section~\ref{sec:methodology}. Any accuracy or F1-score reported for
the Odissi class therefore reflects the classification of a single
instance and does not constitute statistically meaningful evidence of
generalisation for that class. Odissi performance metrics in this study
should be regarded as indicative rather than conclusive. Second, the
pose quality scoring head was trained on dummy labels (constant value
0.5) constructed in the absence of expert-annotated pose quality
ground truth. While the scoring engine functions as a usable
relative quality indicator via the angular deviation computation, the
sigmoid output of the neural network scoring head does not correspond
to a calibrated quality scale grounded in ground-truth expert
annotation. Future work will address this limitation by incorporating
expert-labelled quality annotations and replacing the dummy-label
training regime with supervised quality scoring. The confusion matrix
for the style classification head is shown in Fig.~\ref{fig:confusion}.

\begin{figure}[!t]
\centering
\includegraphics[width=\columnwidth]{figures/confusion_matrix}
\caption{Confusion matrix for the NrityaAI style classification head
on the 8,329-sequence test partition. Rows represent true labels;
columns represent predicted labels. Note the single Odissi test
instance (see Section~\ref{subsec:limitations}).}
\label{fig:confusion}
\end{figure}

Section~V will present the full experimental results and comparative
analysis against prior state-of-the-art systems, including a
per-class breakdown of precision, recall, and F1-score, analysis of
the training convergence behaviour, and a qualitative evaluation of
the pose correction engine on held-out performance clips.

% -----------------------------------------------------------------------
%  BIBLIOGRAPHY
% -----------------------------------------------------------------------
\begin{thebibliography}{10}

\bibitem{openpose2017}
Z. Cao, T. Simon, S.-E. Wei, and Y. Sheikh,
``Realtime multi-person 2D pose estimation using part affinity fields,''
in \textit{Proc. IEEE Conf. Comput. Vis. Pattern Recognit. (CVPR)},
Honolulu, HI, USA, Jul. 2017, pp. 7291--7299.
\url{https://doi.org/10.1109/CVPR.2017.143}

\bibitem{blazepose2020}
V. Bazarevsky, I. Grishchenko, K. Raveendran, T. Zhu, F. Zhang,
and M. Grundmann,
``BlazePose: On-device real-time body pose tracking,''
\textit{arXiv preprint arXiv:2006.10204}, Jun. 2020.

\bibitem{jain2021}
A. Jain, A. Agrawal, M. Dhanasekaran, and B. Raj,
``Recognition of Indian classical dance forms using deep learning,''
\textit{Applied Sciences}, vol. 11, no. 14, p. 6253, Jul. 2021.
\url{https://doi.org/10.3390/app11146253}

\bibitem{challapalli2022}
S. Challapalli and N. Devarakonda,
``Hybrid CNN-PSO-GWO framework for dance pose classification,''
\textit{Knowledge and Information Systems}, vol. 64, no. 9,
pp. 2421--2447, Springer, Sep. 2022.

\bibitem{naik2021}
A. Naik and M. Supriya,
``3D point cloud-based Indian classical dance recognition,''
in \textit{Proc. Int. Conf. Comput. Vis. Bio-Inspired Comput. (ICCVBIC)},
Springer, 2021.

\bibitem{survey2024}
[Author names],
``A survey on computational approaches to Indian classical dance
recognition,''
\textit{Entertainment Computing}, Elsevier, 2024.
\url{https://doi.org/10.1016/j.entcom.2024.100077}

\bibitem{ieee2025}
[Author names],
``Novel deep learning approach for Indian classical dance pose
classification,''
in \textit{IEEE Xplore}, 2025. Article ID: 11113961.

\bibitem{thoutam2022}
V. Thoutam, A. Vyas, B. K. Balabantaray, V. Shruti, V. Shrivastava,
and A. Srivastava,
``Yoga pose estimation and feedback generation using deep learning,''
\textit{Computational Intelligence and Neuroscience},
Wiley-Hindawi, vol. 2022, Article ID 1823795, 2022.
\url{https://doi.org/10.1155/2022/1823795}

\bibitem{swain2022}
S. Swain, A. Mishra, and H. Sharma,
``Real-time yoga monitoring and feedback system using MediaPipe and
CNN+LSTM architecture,''
\textit{Algorithms}, MDPI, vol. 15, no. 11, p. 403, Oct. 2022.
\url{https://doi.org/10.3390/a15110403}

\bibitem{yogavision2024}
[Author names],
``Yoga Vision: CNN+LSTM+MediaPipe fusion for real-time yoga pose
correction,''
in \textit{IEEE Xplore}, 2024. Article ID: 10688319.

\end{thebibliography}

\end{document}
